\SourceLeaf{128}
Fristen z. B. aller 5 oder 10 Jahre, indem man immer die neu hinzugekommenen Erfahrungen mit benutzt, und sieht nur dann den Ueberschuss als theilweis disponibel an, wenn er sich wirklich \textit{vergrössert} hat.

Aber auch abgesehen von diesen beiden Umständen, oder mit andern Worten, auch wenn man einen bestimmten Zinsfuss und eine bestimmte Zahl alljährlich im Durchschnitt beitretender neuer Mitglieder zum Grunde legt, scheint doch die Schwierigkeit der Abschätzung fast unüberwindlich, da die verschiedensten Verhältnisse vom Alter der Ehegatten eintreten, der Wiederverheirathung verwitwet gewordener nicht einmal zu gedenken. In jener Beziehung scheint also der Begutachter nur ungefähr auf Einer Linie zu stehen mit demjenigen, der den Plan von einer der vielen Witwenkassen hätte im Voraus prüfen sollen, die ohne strenge Berücksichtigung des Lebensalters der eintretenden Ehepaare errichtet, fast alle zu Grunde gegangen sind. (Wenn Herr Universitätsrath K. glaubt, dass es auch bei allen diesen Kassen an \textit{Calcül} nicht gefehlt haben werde, so hat er ohne Zweifel Recht; wenn er aber daraus auf die Bodenlosigkeit der \textit{Wahrscheinlichkeitsrechnung} schliessen will, so hat er Unrecht. Allerdings gibt es viele Wörter, mit denen verschiedene Personen verschiedene Bedeutungen verbinden, imgleichen solche, die wissenschaftlich eine sehr bestimmte Bedeutung haben, unter denen man aber im gemeinen Leben oft sehr disparate Dinge zusammenwirft. So ist es mit dem Ausdruck Wahrscheinlichkeitsrechnung bewandt. Im, strengen Sinne verstanden kann von Anwendung derselben in allen den Fällen gar nicht die Rede sein, wo die nöthigen Grundlagen fehlen. Bei allen den gescheiterten Witwenkassen ist bei der Anordnung der Einrichtung von der strengen Wahrscheinlichkeitsrechnung \textit{gar kein} Gebrauch gemacht, sondern nur von vagen Aperçüs. Dies spreche ich hier nur als Thatsache aus, aber nicht als Vorwurf, da in der That eine Basirung auf Wahrscheinlichkeitsrechnung schon darum unmöglich war, weil alle nothwendigen Bedingungen dazu fehlten). Allein in dem vorliegenden Fall ist es zwar unmöglich, ein Endresultat nach der Wahrscheinlichkeitsrechnung aus den einzelnen Elementen zu ermitteln, eben weil diese Elemente fehlen, wohl aber bietet die hundertjährige Erfahrung bei der Kasse selbst, wenn sie auf die rechte Art ausgebeutet wird, einen reichen Schatz zur Grundlage dar. Diese Erfahrungen werden daher erst gesammelt und geordnet werden müssen. Ich setze die Anlegung eines Buches voraus, in welchem von der ersten Stiftung der Gesellschaft an die sämmtlichen Mitglieder, ohne Ausnahme, nach der chronologischen Ordnung des Eintritts verzeichnet werden, nebst allen den Angaben, die für den in Rede stehenden Zweck relevant sind. Allerdings würden diese Erfahrungen ein noch viel fruchtbareres Material darbieten, wenn von sämmtlichen betheiligten Personen auch Geburtsjahr und Tag aufgezeichnet wäre, nemlich von dem eintretenden Professor, von seiner Frau, wenn er schon verheirathet ist, oder, wenn und so oft er sich nach dem Eintritt verehelicht, endlich von den minorennen Kindern, die beim Absterben des Mitgliedes vorhanden sind. Alle diese Dinge aber fehlen, und würden nur eben in Beziehung auf das Mitglied selbst sich noch jetzt in den meisten Fällen ergänzen lassen, aus welchen einzelnen Bestimmungen sich aber wenig oder gar kein Nutzen ziehen liesse. Gleichwohl bleibt das, was sich noch jetzt ohne Zweifel wird zusammenbringen lassen, höchst schätzbar, ich meine nemlich für jedes einzelne Mitglied

1. \textit{Terminus a quo} und \textit{ad quem} der geleisteten Beiträge.

2. \textit{Terminus a quo} und \textit{ad quem} der genossenen Witwenpension in den Fällen wo ein solcher eintrat.

3. \textit{Terminus a quo} und \textit{ad quem} der genossenen Waisenpension, wo nach dem Tode des Vaters oder der Mutter noch minorenne Kinder vorhanden waren.

Die \textit{Grösse} der Geldsumme, die von den Mitgliedern beigetragen, von den Witwen und Waisen erhoben sind, braucht aber gar nicht mit extrahirt werden.

Dies wäre denn das \textit{dritte Requisit} dessen Herbeischaffung, vor Anfang aller Berechnungen, unerlässlich ist. Ich bin mit der Einrichtung des Archivs der Witwenkasse ganz unbekannt, weiss also nicht, ob vielleicht nicht besonders angelegte Bücher, aus denen dieses Material mit Sicherheit, Vollständigkeit
